% Shared preamble for libp2p-raft whitepaper
\usepackage[T1]{fontenc}
\usepackage[utf8]{inputenc}
\usepackage{lmodern}
\usepackage[english]{babel}
\usepackage[
  top=2.6cm,
  bottom=2.6cm,
  left=2.8cm,
  right=2.8cm,
  headheight=14pt
]{geometry}
\usepackage{microtype}
\usepackage{graphicx}
\usepackage{booktabs}
\usepackage{tabularx}
\usepackage{longtable}
\usepackage{array}
\usepackage{enumitem}
\usepackage{hyperref}
\usepackage{xcolor}
\usepackage{listings}
\usepackage{fancyhdr}
\usepackage{titlesec}
\usepackage{tcolorbox}
\tcbuselibrary{skins,breakable}
\usepackage{float}
\usepackage{tikz}
\usetikzlibrary{arrows.meta,decorations.pathreplacing}

\definecolor{brand}{RGB}{0,82,147}
\definecolor{muted}{RGB}{96,96,96}
\definecolor{codebg}{RGB}{245,245,248}

\hypersetup{
  colorlinks=true,
  linkcolor=brand,
  citecolor=brand,
  urlcolor=brand,
  pdftitle={libp2p-raft Architecture Whitepaper},
}

\pagestyle{fancy}
\fancyhf{}
\fancyhead[L]{\small\textcolor{muted}{libp2p-raft --- Architecture Whitepaper}}
\fancyhead[R]{\small\textcolor{muted}{\thepage}}
\renewcommand{\headrulewidth}{0.4pt}

\titleformat{\section}{\large\bfseries\color{brand}}{\thesection}{1em}{}
\titleformat{\subsection}{\normalsize\bfseries}{\thesubsection}{1em}{}
\titleformat{\subsubsection}{\normalsize\bfseries\itshape}{\thesubsubsection}{1em}{}

\lstset{
  basicstyle=\ttfamily\footnotesize,
  breaklines=true,
  columns=fullflexible,
  frame=single,
  rulecolor=\color{black!15},
  backgroundcolor=\color{codebg},
}

\newtcolorbox{invariantbox}[1][]{
  enhanced,
  breakable,
  colback=brand!5,
  colframe=brand!55,
  fonttitle=\bfseries,
  title=#1,
  boxrule=0.6pt,
  arc=2pt,
}

\newtcolorbox{boundarybox}{
  enhanced,
  breakable,
  colback=black!3,
  colframe=black!25,
  title=\textbf{Layer boundary rule},
  boxrule=0.5pt,
}

% Inline code: allow a line break after each underscore so long identifiers
% (e.g. handle_pending_outbound_connection) never overflow the right margin.
\newcommand{\code}[1]{{\def\_{\textunderscore\allowbreak}\texttt{#1}}}

% Sequence diagrams (TikZ)
\tikzset{
  seqactor/.style={draw=brand!60, fill=brand!8, rounded corners=2pt,
    minimum width=2.3cm, minimum height=0.65cm, font=\small\bfseries,
    text height=1.7ex, text depth=0.5ex},
  seqlife/.style={draw=black!30, dashed},
  seqmsg/.style={-{Stealth[length=2mm]}, semithick, draw=black!80},
  seqret/.style={-{Stealth[length=2mm]}, semithick, dashed, draw=black!55},
  seqlbl/.style={above, font=\scriptsize, align=center, fill=white, inner sep=1.2pt},
  seqnote/.style={draw=black!25, fill=codebg, font=\scriptsize, align=center,
    inner sep=2.5pt, rounded corners=1pt},
}
% \seqactors{y_end}{x1/Label1, x2/Label2, ...}: header boxes + lifelines
\newcommand{\seqactors}[2]{%
  \foreach \x/\t in {#2} {%
    \draw[seqlife] (\x,0) -- (\x,#1);
    \node[seqactor] at (\x,0) {\t};
  }}
% \seqmsg[style]{x_from}{x_to}{y}{label}
\newcommand{\seqmsg}[5][seqmsg]{\draw[#1] (#2,#4) -- node[seqlbl]{#5} (#3,#4);}
% \seqnote{x}{y}{text}: local processing on a lifeline
\newcommand{\seqnote}[3]{\node[seqnote] at (#1,#2) {#3};}

% Table columns: ragged-right paragraph cell of a given width
\newcolumntype{L}[1]{>{\raggedright\arraybackslash}p{#1}}
% Width of the last column so a 2-column table fills exactly \linewidth
\newcommand{\restwidth}[1]{\dimexpr\linewidth-#1-2\tabcolsep\relax}

% Block paragraphs (no first-line indent) and avoid overfull lines from \texttt
\setlength{\parindent}{0pt}
\setlength{\parskip}{0.5em plus 0.2em minus 0.1em}
\setlength{\emergencystretch}{3em}
\frenchspacing
% Lists are code-heavy: ragged-right avoids stretched gaps and overflow
\setlist{before=\raggedright}

% Full-width abstract, aligned with body text
\renewenvironment{abstract}{%
  \begin{center}\bfseries\abstractname\end{center}\vspace{-0.5em}\noindent\ignorespaces
}{\par\medskip}
